% Source-derived chapter 12: Computation of $p$-adic periods
% Original source: standard-conjectures-abelian-fourfolds
\chapter{Computation of $p$-adic periods}
\label{standard-conjectures-abelian-fourfolds-chapter-12}
\source{standard-conjectures-abelian-fourfolds}

\begin{remark}[Source section]
\label{standard-conjectures-abelian-fourfolds-chapter-12-source}
\source{standard-conjectures-abelian-fourfolds}
\source{standard-conjectures-abelian-fourfolds}
This chapter reproduces the corresponding section of the retrieved
LaTeX source, including its definitions, statements, examples, and
proofs. It has not yet been translated into Lean.
\notready
\end{remark}

% The source section title was: Computation of $p$-adic periods
\label{computation}
\source{standard-conjectures-abelian-fourfolds}
The aim of this section is to exhibit explicit elements of the space of periods $P_i(F)$ introduced in Definition \ref{defpi} when $F$ is unramified and in Definition \ref{defpiram} when $F$ is ramified. This is done with Propositions  \ref{warmbrinon} and \ref{warmbrinonram} in mind. The construction of these periods is inspired by \cite[\S 9]{Colmez}.


 We will still divide the analysis in two cases, depending  whether $F$ is unramified or not.  Moreover the ramified case will need some extra special cases for $p=2$.

 \subsection*{Unramified case}

In  this subsection we work under   Assumption  \ref{ipos} and we assume moreover that the field $F$  constructed in  Proposition \ref{propendo} is unramified over $\Q_p$.
This means that we have the inclusions \[F \subset  B_{\crys} \subset B_{\dR}\] and that the absolute Frobenius $\varphi$ of $B_\crys$  restricted to $F$ is the  non-trivial element of $\Gal(F/\Q_p)$.


\begin{prop}
\source{standard-conjectures-abelian-fourfolds}
\notready\label{lubintate} (Lubin--Tate)
\source{standard-conjectures-abelian-fourfolds}
There is an element $t_2 \in B_\crys$ which verifies the following properties:
\begin{enumerate}
\item $t_2 \in \Fil ^{1} B_\dR-\Fil ^{2} B_\dR$.
\item  $ \varphi(t_2) \in \Fil ^{0} B_\dR-\Fil ^{1} B_\dR$.
\item $\varphi^2(t_2) = p \cdot t_2$.
\item Both $t_2 $ and  $ \varphi(t_2) $ are invertible in $B_{\crys}$.
\end{enumerate}
\end{prop}
\begin{proof}
This period will appear as the period of an explicit filtered $\varphi$-module. Consider the vector space  $N=\Q_p^2$ with basis $e_1,e_2$ endowed with the Frobenius
\begin{equation}\label{darienzo}  
\source{standard-conjectures-abelian-fourfolds}
\varphi =  \begin{pmatrix} 0 & 1/p \\
1 & 0
 \end{pmatrix}.
 \end{equation}
 Define the filtration to be $\Fil^{-1} N=N, \, \Fil^0 N= \Q_p e_2, \, \Fil^{1} N=0.$ This filtered $\varphi$-module is admissible. To check so notice that there are no non trivial submodules of $N$ stable by $\varphi$ hence the admissibility condition has to be checked only on $N$. 
 Moreover the action of $\varphi$ on $\det N$, the maximal exterior power of $N$, is by multiplication by $-1/p$, which has valuation equal to $-1$, and the only non trivial step of the induced filtration on $\det N$ is at $-1$ as well.
 
 As $N$ is an admissible $\varphi$-module, by Theorem \ref{tsuji} there exists a unique $\Gal_{\Q_p}$-representation $V$ such that $N=D(V)$. 
 Consider the identification $V \otimes B_\crys= N \otimes B_\crys$ given by  Theorem \ref{tsuji}. It allows to write an element of $V$ with respect to the basis $e_1,e_2$. Let $\begin{pmatrix} \alpha \\
\beta
 \end{pmatrix}$ be a fixed nonzero vector in $V$. As the Frobenius must act trivially on it we get from (\ref{darienzo})  the relations
 \[ \alpha= 1/p \cdot \varphi(\beta),\,
    \beta= \varphi(\alpha),\]
  which give the equality
  \[ p \cdot \alpha=  \varphi^2(\alpha).\]
  We claim that $t_2= \alpha$ and $\varphi(t_2)= \beta$ verify all the properties of the proposition.
  
  By Theorem \ref{tsuji} the vector $\begin{pmatrix} \alpha \\
\beta
 \end{pmatrix}$ belongs to $\Fil^0[M \otimes B_\dR]$. This means that it can be written as $\gamma n + \delta e_2$ with $\gamma \in \Fil^{1}B_\dR, \delta \in \Fil^0 B_\dR$ and $n \in N$, which implies that $\alpha \in \Fil^{1}B_\dR$ and  $\beta \in \Fil^{0}B_\dR$. 
  
  Now notice that, for all $f\in F$, the vector $\begin{pmatrix} f\cdot \alpha \\
\varphi(f)\beta
 \end{pmatrix}$ is in the vector space $V=[N\otimes B_\crys]^{\varphi=\id} \cap  \Fil^0[N \otimes B_\dR] $. Hence,  the one dimensional $\Q_p$-vector space of $p$-adic periods relating $\det N$ to $\det V$ is generated by an element in $B_\crys$ of the form $\alpha \beta (f-\sigma( f) )$, for any fixed $f \in F- \Q_p$. As $\det N\otimes B_\crys$ and $\det V\otimes B_\crys$ are free $B_\crys$-modules of rank one the element $\alpha \beta (f-\sigma( f) )$ must be invertible in $B_\crys$, hence $\alpha$ and  $\beta$ are invertible as well.
 
Finally, let us consider the  valuation $\nu$ from Theorem \ref{javier} associated to the filtration. We have showed the inequalities $\nu ( \alpha) \geq 1$ and  $\nu (  \beta)\geq 0$.
On the other hand, as the only non trivial step of the filtration on $\det M$ is  $-1$, we deduce that $\nu ( \alpha \beta (f-\sigma( f) ))=1$ and hence $\nu ( \alpha \beta)=1$. This implies  that the above inequalities are actually equalities and concludes the proof.
 \end{proof}
\begin{corollary}
\source{standard-conjectures-abelian-fourfolds}
\notready\label{brinon}
\source{standard-conjectures-abelian-fourfolds}
Fix a positive integer $i$ and define $\lambda_i \in B_\crys$ as
\[\lambda_i = (t_2/\varphi(t_2))^i .\]
Then $\lambda_i$ is a nonzero element of  $P_i$ and moreover it verifies
\[\lambda_i \cdot\varphi(\lambda_i) = 1/p^i.\] 
\end{corollary}
\begin{proof}
All properties follow from Proposition \ref{lubintate} by direct calculation.
\end{proof}


\subsection*{Ramified case}\label{ramcase}
\source{standard-conjectures-abelian-fourfolds}


In this subsection we work under the Assumption  \ref{ipos} and we assume that the field $F$  constructed in  Proposition \ref{propendo} is ramified over $\Q_p$. This subsection is written in  analogy to the unramified case, although  some extra computations are needed to deal with the case $p=2$.

Define $\sigma$ to be the  non-trivial element of $\Gal(F/\Q_p)$ and $  B_{\crys,F} \subset B_\dR$ as the smallest subring containing $F$ and  $B_{\crys}$.
The inclusions   $  B_{\crys} \subset B_{\crys,F}$ and  $  F \subset B_{\crys,F}$ induce an identification
\[F\otimes_{\Q_p} B_{\crys}=B_{\crys,F}. \]
This allows to extend $\sigma$ and $\varphi$ to two endomorphisms of $\Q_p$-algebras
\[   \sigma, \varphi : B_{\crys,F}    \longrightarrow B_{\crys,F} \]
which commute.

Finally, we will denote by $\Q_{p^n}\subset B_{\crys}$ the unique unramified field extension of $\Q_p$ of degree $n$. 


\begin{prop}
\source{standard-conjectures-abelian-fourfolds}
\notready\label{lubintateram} (Colmez)
\source{standard-conjectures-abelian-fourfolds}
For each uniformizer $\pi \in F$ there exists an element $t_\pi \in B_{\crys,F}$  which verifies the following properties.
\begin{enumerate}
\item $t_\pi \in \Fil ^{1} B_\dR-\Fil ^{2} B_\dR$.
\item $\sigma (t_\pi) \in \Fil ^{0} B_\dR-\Fil ^{1} B_\dR$.
\item  $ \varphi(t_\pi) = \pi  \cdot t_\pi$.
\item Both $t_\pi $ and  $ \sigma(t_\pi) $ are invertible in $B_{\crys,F}$.
\end{enumerate}
\end{prop}
\begin{proof}
This period will appear as the period of an explicit filtered $\varphi$-module. Let $x^2+ax+b$ be the minimal polynomial of $1/\pi$ over $\Q_p$. Consider the vector space  $N=\Q_p^2$ with basis $e_1,e_2$ endowed with the Frobenius
\[ 
\varphi =  \begin{pmatrix} 0 & -b \\
1 & -a
 \end{pmatrix}.
 \]
 Fix an eigenvector $m \in M \otimes F$ satisfying
 \begin{equation}\label{disarli2}  
\source{standard-conjectures-abelian-fourfolds}
\varphi (m)=  \frac{1}{\pi} m
 \end{equation}
 and hence 
  $
\varphi (\sigma(m))=  \frac{1}{\sigma(\pi)} \sigma(m). 
 $
 Define the filtration to be \[\Fil^{-1} (N\otimes F)=N\otimes F, \, \Fil^0 (N\otimes F)= F\cdot ( \sigma(v)), \, \Fil^{1} (N\otimes F)=0.\] We claim that this filtered $\varphi$-module is admissible. To check so notice that there are no non trivial submodules of $N$ stable by $\varphi$ hence the admissibility condition has to be checked only on $N$. Moreover the action of $\varphi$ on $\det N$, the maximal exterior power of $N$, is by multiplication by $b$, which has valuation equal to $-1$, as $\pi$ is a uniformizer of $F$, and the only non trivial step of the induced filtration on $\det N$ is at $-1$ as well.
 
  As $N$ is an admissible filtered $\varphi$-module we can apply Theorem \ref{tsuji} and consider the $\Gal_{\Q_p}$-representation $V$ which verifies  $N=D(V)$. Now notice that the action of $\varphi$ on $N$ respects the filtration, this means that the field $F$ acts on $N$ as filtered $\varphi$-module. Then Theorem \ref{tsuji}  implies that
 $F$ acts on $V$ as well and that the canonical isomorphism  $V \otimes B_\crys= N \otimes B_\crys$ is $F$-equivariant. In particular we can find an eigenvector $v \in V \otimes F$ and a scalar $\alpha \in B_{\crys,F}$ such that 
  \begin{equation}\label{disarli4}  
\source{standard-conjectures-abelian-fourfolds}
\alpha \cdot m=   v
 \end{equation}
 and hence $ \sigma(\alpha) \cdot \sigma(m)=   \sigma(v)$.
 
As the Frobenius acts trivially on $V$, relations (\ref{disarli2}) and (\ref{disarli4}) imply $ \varphi(\alpha) = \pi  \cdot \alpha$. We claim now that the period $t_\pi=\alpha$ satisfies all the properties of the statement. Indeed, as $v$ and $m$ generates the same free  $B_{\crys,F}$-module of rank one, the scalar $\alpha$ must be invertible in $B_{\crys,F}$, and the same argument applies for $ \sigma(\alpha)$. As the filtration is trivial on $V \otimes F$, relation   (\ref{disarli4}) implies that $\alpha \in \Fil ^{1} B_\dR-\Fil ^{2} B_\dR$ and again the same argument applies for $ \sigma(\alpha)$.
 \end{proof}
\begin{corollary}
\source{standard-conjectures-abelian-fourfolds}
\notready\label{brinonram}
\source{standard-conjectures-abelian-fourfolds}
Suppose that there is a uniformizer $\pi\in F$ such that 
\[\sigma(\pi)=-\pi.\]
Let $\sqrt{a}\in \Q_{p^2}$ be an element of the quadratic unramified extension of $\Q_p$ such that
\[\varphi(\sqrt{a}) = -\sqrt{a}.\]
Define $\lambda_i \in B_{\crys,F}$ as
\[\lambda_i = (\sqrt{a}  \cdot t_\pi/ \sigma(t_\pi))^i .\]
Then $\lambda_i$ is a nonzero element of  $P_i(F)$ and moreover it verifies
\[\lambda_i \cdot\sigma(\lambda_i) = a^i.\] 
\end{corollary}
\begin{proof}
All properties follow from  Proposition \ref{lubintateram} by direct calculation.
\end{proof}

\begin{remark}
\source{standard-conjectures-abelian-fourfolds}
\notready\label{remtame}
\source{standard-conjectures-abelian-fourfolds}
Let us recall the list of quadratic extensions of $\Q_p$. If $p\neq 2$, and if $a$ is an integer which is not a square modulo $p$, then $\Q_p(\sqrt{a})$ is the unramified quadratic extension, whereas $\Q_p(\sqrt{p})$ and $\Q_p(\sqrt{ap})$ are the two ramified quadratic extensions. 
If $p=2$,  $\Q_2(\sqrt{5})$ is the unramified quadratic extension, whereas $\Q_2(\sqrt{2}),\Q_2(\sqrt{6}),\Q_2(\sqrt{10}),\Q_2(\sqrt{14}),\Q_2(\sqrt{3})$ and $\mathbb{Q}_2(\sqrt{-1})$ are the six ramified quadratic extensions. 

In particular, it is almost always possible to find a uniformizer $\pi\in F$ such that 
$\sigma(\pi)=-\pi$ as in  Corollary \ref{brinonram}. The only exceptions are for  $p=2$ and  $F=\mathbb{Q}_2(\sqrt{3})$ or  $F=\mathbb{Q}_2(\sqrt{-1})$.  What follows treats these two exceptions.
\end{remark}




\begin{corollary}
\source{standard-conjectures-abelian-fourfolds}
\notready\label{brinonwild}
\source{standard-conjectures-abelian-fourfolds}
We keep notation from Proposition  \ref{lubintateram} and we work  with $p=2$ and the field  $F=\mathbb{Q}_2(\sqrt{-1})$.  Let $\alpha \in \Q_{2^4}(\sqrt{-1})$ be an element  such that
\[\varphi(\alpha) =  \sqrt{-1} \cdot \alpha.\]
Define $\lambda_i \in B_{\crys,F}$ as
\[\lambda_i = (\alpha \cdot t_{1-\sqrt{-1}} /\sigma(t_{1-\sqrt{-1}}))^i .\]
Then $\lambda_i$ is a nonzero element of  $P_{i}(F) $ and  moreover it verifies
\[\lambda_i \cdot\sigma(\lambda_i) = (\alpha  \cdot \sigma(\alpha) )^i.\] 
\end{corollary}
\begin{proof}
All properties follow from Proposition \ref{lubintateram}  by direct calculation.
\end{proof}



\begin{prop}
\source{standard-conjectures-abelian-fourfolds}
\notready\label{idearef} (Colmez)
\source{standard-conjectures-abelian-fourfolds}
For each uniformizer $\pi \in F$ there exists an element $t_{2,\pi} \in B_{\crys,F}$  which verifies the following properties.
\begin{enumerate}
\item $t_{2,\pi}  \in \Fil ^{1} B_\dR-\Fil ^{2} B_\dR$.
\item $\sigma (t_{2,\pi} ) , \varphi (t_{2,\pi} ) , \varphi \circ \sigma (t_{2,\pi} ) \in \Fil ^{0} B_\dR-\Fil ^{1} B_\dR$.
\item  $ \varphi^2(t_{2,\pi} ) = \pi \cdot t_{2,\pi} $.
\item The elements $t_{2,\pi}  ,  \sigma(t_{2,\pi} ), \varphi (t_{2,\pi} )$ and $ \varphi \circ \sigma (t_{2,\pi} )  $ are invertible in $B_{\crys,F}$.
\end{enumerate}
\end{prop}
\begin{proof}
This period will appear as the period of an explicit filtered $\varphi$-module. Let $x^2+ax+b$ be the minimal polynomial of $1/\pi$ over $\Q_p$.  Define the $2 \times 2$ matrix
\[ 
C =  \begin{pmatrix} 0 & -b \\
1 & -a
 \end{pmatrix}.
 \]
 Consider the vector space  $N=\Q_p^4$  endowed with the Frobenius given by the block matrix
 \[ 
\varphi =  \begin{pmatrix} 0_2 & C \\
I_2 & 0_2
 \end{pmatrix},
 \]
where $I_2$ and $0_2$ are the identity and the zero $2 \times 2$  matrices.  

 Notice  that we have \[ 
\varphi^2 =  \begin{pmatrix} C & 0_2 \\
0_2 & C
 \end{pmatrix}.
 \]
 Hence the action of $\varphi^2$ induces a decomposition of eigenspaces
  \[N \otimes F=N_{1/\pi}\otimes N_{1/\sigma(\pi)}.\]
 Fix a nonzero eigenvector $m \in N_{1/\pi}$ hence satisfying
 \begin{equation}\label{troilo}  
\source{standard-conjectures-abelian-fourfolds}
\varphi^2 (m)=  \frac{1}{\pi} m.
 \end{equation}
Define the filtration to be \[\Fil^{-1} (N\otimes F)=N\otimes F, \, \Fil^0 (N\otimes F)= N_{1/\sigma(\pi)}\oplus F\cdot \varphi (m ) , \, \Fil^{1} (N\otimes F)=0.\] 
We claim that this filtered $\varphi$-module is admissible. First notice that there are no non trivial submodules of $N$ stable by $\varphi$. Indeed the action of $\varphi$ on $N$ satisfies the equation $x^4+ax^2+b=0$. This polynomial is irreducible over $\Q_p$ as it is also the minimal polynomial of $1/\sqrt{\pi}$, which has valuation $1/4$.
 We deduce that the admissibility condition has to be checked only on $N$. Moreover the action of $\varphi$ on $\det N$, the maximal exterior power of $N$, is by multiplication by $b$, which has valuation equal to $-1$ and the only non trivial step of the induced filtration on $\det N$ is at $-1$ as well.
 
As $N$ is an admissible filtered $\varphi$-module we can apply Theorem \ref{tsuji} and consider the $\Gal_{\Q_p}$-representation $V$ which verifies  $N=D(V)$. 
Now notice that the action of $\varphi^2$ on $N$ respects the filtration, this means that the field $F$ acts on $N$ as filtered $\varphi$-module. 
Then Theorem \ref{tsuji}  implies that
 $F$ acts on $V$ as well and that the canonical isomorphism  $V \otimes B_\crys= N \otimes B_\crys$ is $F$-equivariant. In particular there is a decomposition in eigenspaces  
 \[V \otimes F=V_{1/\pi}\otimes V_{1/\sigma(\pi)}\]
 and, if we fix  a nonzero eigenvector $v \in V_{1/\pi}$, there exist two periods $\alpha, \beta$ in $B_{\crys,F}$ such that 
  \begin{equation}\label{troilo2}  
\source{standard-conjectures-abelian-fourfolds}
\alpha \cdot m+ \beta \cdot \varphi(m)=   v.
 \end{equation}
 
As the Frobenius acts trivially on $V$, relations (\ref{troilo}) and (\ref{troilo2}) imply $ \varphi(\alpha) = \beta$ and $ \varphi^2(\alpha) = \pi \alpha$ . We claim now that the period $t_{2,\pi}=\alpha$ satisfies all the properties of the statement. 

 By Theorem \ref{tsuji} we have $v\in \Fil^0[N_{1/\pi} \otimes B_\dR]$. This means that the vector $v$ can be written as $\gamma n + \delta \varphi(m)$ with $\gamma \in \Fil^{1}B_\dR, \delta \in \Fil^0B_\dR$ and $n \in N_{1/\pi}$, which implies that $\alpha \in \Fil^{1}B_\dR$ and  $\beta \in \Fil^{0}B_\dR$. 

Now notice that, for all $x\in \Q_{p^2}$, the vector $x\alpha \cdot m+ \varphi(x)\varphi(\alpha) \cdot \varphi(m)$ is in the vector space $V=[N\otimes B_\crys]^{\varphi=\id} \cap  \Fil^0[N \otimes B_\dR] $. 
Moreover, it is in the eigenspace $V_{1/\pi}$. 
Hence,  the one dimensional $F$-vector space of $p$-adic periods relating $\det N_{1/\pi}$ to $\det V_{1/\pi}$ is generated by an element in $B_{\crys,F}$ of the form $\alpha \varphi(\alpha) (x-\varphi( x) )$, for any fixed $x \in \Q_{p^2} - \Q_p$. 
As $\det N_{1/\pi}\otimes B_{\crys,F}$ and $\det V_{1/\pi}\otimes B_{\crys,F}$ are free $B_{\crys,F}$-modules of rank one the element $\alpha \varphi(\alpha) (x-\varphi( x) )$ must be invertible in $B_{\crys,F}$, hence $\alpha$ and  $\varphi(\alpha)$ are invertible as well.
 
 Finally, let us  consider the valuation $\nu$ from Theorem \ref{javier} associated  to the filtration. 
 First of all, as the only non trivial step of the filtration on $\det N_{1/\pi}$ is  $-1$, we deduce that $\nu ( \alpha \varphi(\alpha) (f-\sigma( f) ))=1$ and hence $\nu ( \alpha \varphi(\alpha))=1$. 
 On the other hand we have showed the inequalities $\nu ( \alpha) \geq 1$ and  $\nu (  \varphi(\alpha))\geq 0$. 
 Combining them with $\nu ( \alpha \varphi(\alpha)  )=1$ we deduce that these inequalities are actually equalities.
 
 This shows all the desired properties on $\alpha$ and $ \varphi(\alpha)$. 
 To show the analogous properties on $\sigma (\alpha)$ and $ \varphi \circ \sigma (\alpha)$ one applies $\sigma$ on the relation (\ref{troilo2}) and argues as above replacing the eigenspaces $N_{1/\pi}$ and $V_{1/\pi}$ with the eigenspaces $N_{1/\sigma(\pi)}$  and $V_{1/\sigma(\pi)}$.
\end{proof}

\begin{prop}
\source{standard-conjectures-abelian-fourfolds}
\notready\label{lubintatewild2} (Colmez)
\source{standard-conjectures-abelian-fourfolds}
Consider the ring $B_{\crys,F}$ for the prime $p=2$ and $F=\Q_2(\sqrt{3}).$
There  exists an element $t_{f} \in B_{\crys,F}$ which verifies the following properties.
\begin{enumerate}
\item $t_{f} \in \Fil ^{1} B_\dR-\Fil ^{2} B_\dR$.
\item $\sigma (t_{f}) , \varphi (t_{f}) , \varphi \circ \sigma (t_{f}) \in \Fil ^{0} B_\dR-\Fil ^{1} B_\dR$.
\item  $ \varphi^2(t_{f}) = (1-\sqrt{-1})  \cdot t_{f}$.
\item The elements $t_{f}  ,  \sigma(t_{f} ), \varphi (t_{f} )$ and $ \varphi \circ \sigma (t_{f} )  $ are invertible in $B_{\crys,F}$.
\end{enumerate}
\end{prop}
\begin{proof}
This is a special case of Proposition \ref{idearef} up to a little subtlety. Consider first the ramified extension $F=\mathbb{Q}_2(\sqrt{-1})$ and the uniformizer $\pi= 1-\sqrt{-1}. $ Then Proposition  \ref{idearef}  will give a period $t_{2,\pi}$ which verifies all the properties of the statement except that it is constructed in $B_{\crys,\mathbb{Q}_2(\sqrt{-1})}$. On the other hand we have the equality \[B_{\crys,\mathbb{Q}_2(\sqrt{-1})} = B_{\crys,\mathbb{Q}_2(\sqrt{3})}.\]
Notice though that the two descriptions of this ring exchange the role of $ \sigma $ with $ \varphi \circ \sigma$. As  Proposition \ref{idearef}  is stable under this exchange, the period $t_f=t_{2,\pi}$ does satisfy all the desired properties.
\end{proof}
\begin{corollary}
\source{standard-conjectures-abelian-fourfolds}
\notready\label{brinonwild2}
\source{standard-conjectures-abelian-fourfolds}
  Consider $F=\Q_2(\sqrt{3})$ and let $\alpha \in \Q_{2^4}(\sqrt{3})=\Q_{2^4}(\sqrt{-1})$ be an element  such that
\[\varphi(\alpha) =  \sqrt{-1} \cdot \alpha.\]
Define $\lambda_i \in B_{\crys,F}$ as
\[\lambda_i = \left(\frac{\alpha \cdot t_{f} \cdot \varphi(t_{f}) }{\sigma(t_{f})   \cdot (\varphi \circ \sigma) (t_{f})}\right)^i .\]
Then $\lambda_i$ is a nonzero element of  $P_{i}(F)$ and  moreover it verifies
\[\lambda_i \cdot\sigma(\lambda_i) = (\alpha  \cdot \sigma(\alpha))^i.\] 
\end{corollary}
\begin{proof}
All properties follow from Proposition \ref{lubintatewild2}  by direct calculation.
\end{proof}
